\documentclass[10pt,a4paper]{amsart}
\usepackage[margin=19mm]{geometry}
\usepackage{amsmath,amssymb,mathtools}
\usepackage[scr=rsfs,cal=euler]{mathalfa}
\usepackage{xfrac}
\usepackage[unicode,hidelinks,bookmarks=false]{hyperref}
\newcommand{\ie}{i.e.\ }
\newcommand{\cf}{cf.\ }
\newcommand{\resp}{respectively\ }
\newcommand{\Subsection}{\subsection}
\providecommand{\nobreakdash}{\nobreak\hbox{-}}
\setlength{\emergencystretch}{2em}
\multlinegap0pt
\usepackage{mathtools}
\usepackage{amssymb}
\usepackage{bbm}
\usepackage{xcolor}
\usepackage[normalem]{ulem}
\newtheorem{lem}{Lemma}
\newtheorem{cor}{Corollary}
\newtheorem{prop}{Proposition}
\newtheorem{thm}{Theorem}
\newcommand{\CC}{\mathcal{C}}
\newcommand{\CM}{\mathcal{M}}
\def\N{\mathbb N}
\newcommand{\norm}[1]{\left\Vert#1\right\Vert}
\newcommand{\ups}{\upsilon}
\title{Non-Commutative Maximal Inequalities for State Preserving Actions of amenable groups}
\author{Panchugopal Bikram, Hariharan G, Sudipta Kundu,, Diptesh Saha}
\date{Source: arXiv:2512.24751v1 (CC BY 4.0)}
\begin{document}
\maketitle

\noindent\emph{Source and licence.} Non-Commutative Maximal Inequalities for State Preserving Actions of amenable groups. Version arXiv:2512.24751v1 (CC BY 4.0), distributed by arXiv under the Creative Commons Attribution 4.0 International licence, http://creativecommons.org/licenses/by/4.0/, as recorded in the arXiv licence metadata for this exact version and shown on the abstract page https://arxiv.org/abs/2512.24751v1. Full licence notice: \texttt{licenses/arxiv-2512.24751-CC-BY-4.0.txt}.\par\smallskip

\noindent\textit{Editorial scope.} Counted unit: the article's thm principle (source0.tex lines 1684--1692). The article's complete original proof of it is delivered with it (source0.tex lines 1694--1799, one \texttt{proof} environment of 106 lines). No mathematics is written, repaired or re-derived: the statement and the proof are the article's own lines, in the article's own notation, and no retained line is substituted or re-flowed.  Retained uncounted as the source-local context the proof needs: lines 900--922; lines 1094--1113; lines 1671--1673.  Nothing was omitted from any retained span, so the package carries no omission record.  No retained line contains a citation, so the wrapper carries no bibliography and no bibliography entry.  No retained line contains a figure environment, an image-inclusion command or drawing code, so no image is delivered and the package declares no asset dependency.  Editorial formatting only, in addition: an AMS \texttt{amsart} preamble that restates the article's own declarations and macros used by the retained lines, namely the environments \texttt{lem}, \texttt{cor}, \texttt{prop}, \texttt{thm} on separate counters, together with the standard packages those lines need.  The retained environments print in this document as Lemma 1, Corollary 1, Proposition 1, Theorem 1 in order of appearance, whereas the article prints the same items with the numbering of its own full document.  The complete native source of the article as distributed in the arXiv e-print for this exact version is bundled unchanged as \texttt{sources/191925/source0.tex}.
\begin{lem}\label{uniform conv lem}
Let \( \CM_1 \) and \( \CM_2 \) be two von Neumann algebras. Let \( \varphi \) be an f.n.s. weight on \( \CM_1 \), and let \( \rho \) and \( \tau \) be an f.n. state and an f.n. tracial state on \( \CM_2 \), respectively.
Suppose that
\[
\CM_1 \,\otimes\, \CM_2 \subseteq \mathcal{B}\big(L^2(\CM_1 \,\otimes\, \CM_2,\, \varphi \otimes \rho)\big).
\]
Further, let $S$ be an arbitrary set, and $
\{ p_j{(s)} : j\in \mathbb{N},~ s \in S  \} \subseteq \mathcal{P}(\CM_1)$ and  $
\{ x_j{(s)} : j\in \mathbb{N},~ s \in S \} \subseteq \CM_1 \,\otimes\, \CM_2
$ be a  family  of projections and a 
 family of operators, respectively, such that 
$$ 0< \varphi(p_j(s)) < \infty \text{ and } 0\leq x_j(s) \leq p_j(s) \otimes 1, ~ \text{ for all } j \in \mathbb{N}, ~ s \in S.
$$
Define
\[
f_j(s) := \frac{p_j(s) \otimes 1 - x_j(s)}{\varphi(p_j(s))}, \quad \text{for } j \in \mathbb{N}, ~s\in S.
\]
Then the following statements are equivalent:
\begin{itemize}\setlength\itemsep{.7em}
    \item[(1)] \( \displaystyle \lim_{j \to \infty} (\varphi \otimes \rho)(f_j(s)) = 0 \), uniformly in \( s \in S \),
    \item[(2)] \( \displaystyle \lim_{j \to \infty} (\varphi \otimes \tau)(f_j(s)) = 0 \), uniformly in \( s \in S \).
\end{itemize}
\end{lem}

\begin{cor}\label{maximal ineq for averages-cor}


Let $\CM$ be a finite von Neumann algebra with a faithful normal tracial state $\tau$, and let $(\CM, G, \alpha, \rho)$ be a kernel. For all $j \in \N$, let $\phi_j \in \CM_{* s}$ satisfy $\norm{\phi_j} \leq \frac{1}{4^{s_j}}$. Fix $ \ups > 0$. Then for all $n \in \N$, there exist a sequence $\left(\gamma_j(n)\right)_{j \in \N}$ of positive real numbers and projections $\{p^{\ups}_{n,j} : j \in \N\}$ in $\CM$ such that:
	\begin{enumerate}\setlength\itemsep{.7em} 
		\item $ \displaystyle \lim_{j \to \infty} \gamma_j(n) = 0$, uniformly in $n \in \N$.  
		\item $\rho(1 - p^{ \ups}_{n,j}) \leq (1 + \ups)\gamma_j(n) + \ups$.  
		\item For all $j \in \N$,
		\begin{align}\label{max-cond}
		    A_k(\phi_j)(x) \leq  \frac{c s_j}{2^{s_j}} \rho(x), \quad \text{for all } x \in p^{ \ups}_{n,j} \CM_+ p^{\ups}_{n,j} \text{ and } k \in [n].
		\end{align}
	\end{enumerate}
	The constant $c$ depends only on the group $G$.






\end{cor}

\begin{prop}\label{complete wrt bau-functionals}
	Every b.a.u.-Cauchy sequence $(\nu_n)_{n \in \N}$ in $\CM_*$ converges in the b.a.u. topology to some $\nu \in \CM_*$.
\end{prop}

    	\begin{thm}[Banach principle]\label{Banach principle}
		Suppose $G$ is a \textit{l.c.s.c.} amenable group with admissible F{\o}lner sequence $(F_{n})_{n \in \mathbb{N}}$, and let $\CM$ be a finite von Neumann algebra equipped with a faithful normal tracial state $\tau$ and a faithful normal state $\rho$. 
		
		Let $(\CM, G, \alpha)$ be a non-commutative dynamical system such that $\rho \circ \alpha_g = \rho$ for all $g \in G$, i.e., $(\CM, G, \alpha, \rho)$ forms a kernel. Then, the following set 
		\begin{align*}
			\CC = \{ \phi \in \CM_{*s} \mid A_n(\phi) \text{ converges in b.a.u.} \}
		\end{align*}
		is closed in $\CM_{*s}$.
	\end{thm}

    	\begin{proof}
		Let $(\phi_n)$ be a sequence in $\mathcal{C}$, converging to $\phi \in \mathcal{M}_{*s}$ in norm, i.e.,
		$$ \lim_{n \to \infty} \| \phi_n - \phi \| = 0. $$
		Then, we need to show that $\phi \in \mathcal{C}$. 
		
		First, we note that since $\rho \in \mathcal{M}_{*+}$, there exists a unique $X \in L^1(\mathcal{M}, \tau)_+$ such that $\rho(x) = \tau(Xx)$ for all $x \in \mathcal{M}$. Then, for any $s > 0$, consider the projection $q_s := 1_{(1/s, s)}(X) \in \mathcal{M}$. Observe that $(1 - q_s) \xrightarrow{s \to \infty} 0$ in the strong operator topology (SOT). 
		
		Let $\theta > 0$ be arbitrary. Therefore, there exists an $s_0 > 0$ such that $\tau(1 - q_{s_0}) < \frac{\theta}{2}$. Further, this implies $X q_{s_0} \leq s_0 q_{s_0}$. Thus, for all $0 \neq x \in q_{s_0} \mathcal{M}_+ q_{s_0}$, we have 
		\begin{align}\label{rho-tau}
			\begin{split}
				\frac{\rho(x)}{\tau(x)} = \frac{\tau(Xx)}{\tau(x)}&= \frac{\tau(X q_{s_0} x)}{\tau(x)} \quad \text{ (since $q_{s_0} x = x$)}\\
				&\leq \frac{\tau(s_0 x)}{\tau(x)} = s_0.
			\end{split} 
		\end{align}
		
		Since $ \lim_{n \to \infty} \| \phi - \phi_n \| = 0 $, possibly passing through a subsequence, we can assume that 
		\begin{align*}\label{coneq-3}
			\|\phi - \phi_j\| < \frac{1}{4^{s_j}} \text{ for all } j \in \mathbb{N}.
		\end{align*}
		
		Furthermore, note that $\phi - \phi_j \in \mathcal{M}_{*s}$ for all $j \in \mathbb{N}$. Now, using Corollary \ref{maximal ineq for averages-cor} for $n = N$, $\upsilon = \frac{1}{2^i}$, we obtain positive reals $({\gamma_j(N)})_{N \in \mathbb{N}}$ and projections $ q_{N,j}^{i} := q^{ \frac{1}{2^i}}_{ N, j} $ in $\mathcal{M}$ for $ j , N \in \mathbb{N}$ such that:
		\begin{enumerate}\setlength\itemsep{.7em}
			\item $ \displaystyle \lim_{j \to \infty} \gamma_j(N) = 0$, uniformly in $N \in \mathbb{N}$,   
			\item $ \rho(1- q^{i}_{N,j}) \leq ( 1+ \frac{1}{2^i})\gamma_j(N) + \frac{1}{2^i} $, and
			\item for all $j \in \mathbb{N}$, 
			\[
			A_n(\phi - \phi_j)(x) \leq  \frac{cs_j}{2^{s_j}}  \rho(x) \text{ for all } x \in q^{i}_{N,j} \mathcal{M}_+ q^{i}_{N, j} \text{ and } n \in [N].
			\]
		\end{enumerate}
        Now denote $q_{N,j} = q^{j}_{N,j}$ \\
        We have,
        \begin{align*}
        \rho(1-q_{N,j}) \leq (1+ \frac{1}{2^j})\gamma_j(N) + \frac{1}{2^j}
        \end{align*} for all $N \in \mathbb{N}$.
		Now, it immediately follows that
		\[
		\displaystyle \lim_{j \to \infty}  \rho(1- q_{N, j}) = 0, \text{ uniformly in }  N \in \mathbb{N}. 
		\]
		We use Lemma \ref{uniform conv lem}. Let $\CM_{1}=L^{\infty}(G)$ and $\CM_{2}=\CM$, take $S$ to be the set of natural numbers $\mathbb{N}$, let  $x_{j}(n)=q_{j,n} \text{ and } p_{j}(n)=\chi_{F_{n}}  \forall n \in \mathbb{N}$. Note that $f_{j}(N)=\frac{p_{j}(N) \otimes 1-q_{j,N}}{\mu(F_{N})},\forall N \in \mathbb{N}$. We have,
        \begin{align*}
        &\lim_{j \rightarrow \infty}\frac{(\varphi \otimes \rho)(\chi_{F_{N}}\otimes(1-q_{N,j}))}{\varphi(\chi_{F_{N}})} = 0 \\
        \implies & \lim_{j \rightarrow \infty}\frac{(\varphi \otimes \tau)(\chi_{F_{N}}\otimes(1-q_{N,j}))}{\varphi(\chi_{F_{N})}} = 0 \text{ by Lemma 3.10 } \\
        \implies & \lim_{j \rightarrow \infty}\frac{\varphi(\chi_{F_{N}})\tau(1-q_{N,j}))}{\varphi(\chi_{F_{N}})} = 0 \\
        \implies & \lim_{j \rightarrow \infty}\tau(1-q_{N,j})=0, \text{ uniformly in N}\in \mathbb{N}.
        \end{align*}
		
		
		\noindent Therefore, there exists a subsequence $(j_l)_{l \in \mathbb{N}}$ such that  
		\begin{align*}
			\tau(1 - q_{N, j_l}) < \frac{\theta}{2^{l+2}},  \quad \text{for all } N \in \mathbb{N}.
		\end{align*}
		
		
		Since $\phi_j \in \mathcal{C}$ for all $j \in \mathbb{N}$, there exists a sequence of projections $(f_j)$ in $\mathcal{P}(\mathcal{M})$ and a strictly increasing sequence $(N_j) \subseteq \mathbb{N}$ such that  
		\begin{align*}
			\begin{split}
				(1)\quad &\tau(1 - f_j)  < \frac{\theta}{2^{j+2}},  \quad \text{for all } j \in \mathbb{N}, \\
				(2)\quad &\sup_{x \in f_j M_+ f_j, x \neq 0} \frac{\left| \big(A_n(\phi_j) - A_m(\phi_j) \big)(x) \right|}{\tau(x)} < \frac{1}{2^j},  \quad \text{for all } n, m \geq N_j.
			\end{split}
		\end{align*}
		
		\noindent Now, for all $l \in \mathbb{N}$, we define $q_l := q_{N_{j_{l+1}}, j_l}$. Then, we note that  
		\begin{align*}
			\begin{split}
				(1)\quad &\tau(1 - q_l) \leq \frac{\theta}{2^{l+2}}, \\
				(2)\quad &A_n(\phi - \phi_{j_l}) (q_l x q_l) \leq   \frac{cs_{j_{l}}}{2^{s_{j_l}}} \rho(q_l x q_l),  
				\quad \text{for all }n \in [N_{j_{l +1}}].
			\end{split}
		\end{align*}
		
		
		
		\noindent Moreover, for the subsequence $(j_l)$, we have  
		\begin{align*}
			\tau(1 - f_{j_l}) \leq \frac{\theta}{2^{j_l+2}}.
		\end{align*}
		
		\noindent Finally, consider the projection  
		\begin{align*}
			e= \left( \underset{l \in \mathbb{N}}{\bigwedge} q_l \right) \bigwedge \left( \underset{l \in \mathbb{N}}{\bigwedge} f_{j_l} \right) \bigwedge q_{s_0}.
		\end{align*}
		Then, we have $\tau(1 - e) < \theta$ and, for all $0 \neq x \in e \mathcal{M}_+ e$, assuming $N_{j_l} \leq n, m \leq N_{j_{l+1}}$, we obtain  
		\begin{align*}
			&\frac{\left| \big(A_n( \phi) - A_m(\phi) \big) (x) \right|}{\tau(x)}\\
			=&  \frac{\left| \big(A_n( \phi) - A_n(\phi_{j_l}) + A_n(\phi_{j_l}) - A_m(\phi_{j_l}) + A_m(\phi_{j_l}) - A_m(\phi) \big) (x) \right|}{\tau(x)}\\
			\leq &  \frac{\left| \big(A_n( \phi) - A_n(\phi_{j_l}) \big)(x)  \right| }{\tau(x) }
			+ \frac{\left| \big(A_n(\phi_{j_l}) - A_m(\phi_{j_l}) \big) (x) \right|}{\tau(x)} 
			+ \frac{\left| \big( A_m(\phi_{j_l}) - A_m(\phi) \big) (x) \right|}{\tau(x)}\\
			\leq &  \frac{\left| \big(A_n( \phi - \phi_{j_l}) \big)(x)  \right| }{\tau(x) }
			+ \frac{\left| \big(A_n(\phi_{j_l}) - A_m(\phi_{j_l}) \big) (x) \right|}{\tau(x)} 
			+ \frac{\left| \big( A_m(\phi_{j_l} - \phi) \big) (x) \right|}{\tau(x)}\\
			\leq &  \left(\frac{\left| \big(A_n( \phi - \phi_{j_l}) \big)(x)  \right| }{\rho(x) } 
			+ \frac{\left| \big( A_m(\phi_{j_l} - \phi) \big) (x) \right|}{\rho(x)} \right) 
			\frac{\rho(x)}{\tau(x)} 
			+ \frac{\left| \big(A_n(\phi_{j_l}) - A_m(\phi_{j_l}) \big) (x) \right|}{\tau(x)}\\
			\leq&  \left(\frac{2cs_{j_{l}}}{2^{s_{j_{l}}}}\right)s_{0} + \frac{1}{2^{s_{j_{l}}}} = \frac{(2cs_{j_{l}}s_{0} +1)} {2^{s_{j_{l}}}}, 
			\quad (\text{by Eq. } \ref{rho-tau}). 
		\end{align*}	
		
		\noindent Since $j_l \to \infty$ as $m,n \to \infty$, it follows that  
		\begin{align*}
			\lim_{m,n \to \infty} \sup_{x \in e M_+ e, x \neq 0} \frac{\left| \big(A_n( \phi) - A_m(\phi) \big) (x) \right|}{\tau(x)} = 0.
		\end{align*}
		
		\noindent Therefore, $(A_n(\phi))_{n \in \mathbb{N}}$ is Cauchy in b.a.u. Thus, by Proposition \ref{complete wrt bau-functionals}, it follows that $(A_n(\phi))_{n \in \mathbb{N}}$ is convergent in b.a.u., which implies that the set $\mathcal{C}$ is closed in $\mathcal{M}_{*s}$.		
	\end{proof}

\end{document}
